(1) 由 \(m\dot{\boldsymbol v}=q(\boldsymbol E+\boldsymbol v\times\boldsymbol B)\)，在 \(xy\) 面写成 \(\dot v_x=qE/m+\omega_c v_y\)、\(\dot v_y=-\omega_c v_x\)，其中 \(\omega_c=qB/m\)。符号来自 \(\boldsymbol v\times B\hat z=(Bv_y,-Bv_x,0)\)，正电荷回旋为顺时针。
(2) 恒定速度须满足 \(\boldsymbol E+\boldsymbol v_d\times\boldsymbol B=0\)，故 \(\boldsymbol v_d=(\boldsymbol E\times\boldsymbol B)/B^2=-(E/B)\hat y\)。此漂移与粒子的正质量、电量无关，电荷符号只决定回旋方向，不改变均匀场的 \(E\times B\) 漂移方向。
(3) 令 \(u=v_x\)、\(w=v_y+E/B\)，则 \(\dot u=\omega_cw\)、\(\dot w=-\omega_cu\)；初值 \(u(0)=0,w(0)=E/B\)。所以 \(v_x=(E/B)\sin\omega_ct\)、\(v_y=(E/B)(\cos\omega_ct-1)\)。积分得 \(x(t)=E(1-\cos\omega_ct)/(B\omega_c)\)、\(y(t)=E\sin\omega_ct/(B\omega_c)-Et/B\)，为在匀速漂移上叠加闭合回旋的旋轮轨迹。
(4) 回旋周期 \(T_c=2\pi/\omega_c\) 后 \(\Delta x=0\)、\(\Delta y=-ET_c/B\)，平均速度即 \(-E\hat y/B\)。电场一周期净功为 \(qE\Delta x=0\)；瞬时动能虽振荡，却在每周期复原。若碰撞、非均匀场、辐射反作用显著，或 \(E/B\) 接近光速，则必须重新分析，不能将非相对论漂移无限外推。
